\small
\arrayrulecolor{black!22}
\begin{longtable}{L{0.14\linewidth} L{0.19\linewidth} L{0.18\linewidth} L{0.31\linewidth} L{0.16\linewidth}}
\hline
\textbf{Workflow} & \textbf{Trigger} & \textbf{Required input} & \textbf{Required updates} & \textbf{Exit rule}\\
\hline
\endfirsthead
\hline
\textbf{Workflow} & \textbf{Trigger} & \textbf{Required input} & \textbf{Required updates} & \textbf{Exit rule}\\
\hline
\endhead
Global no-go guardrail inheritance & Every phase preparation, implementation prompt, code review, object documentation update, closeout, pruning decision, or generated-view refresh. & top-level global\_no\_go\_guardrails;\newline active phase scratchpad;\newline proposed module/\allowbreak{}declaration/\allowbreak{}object changes;\newline operational vs mirror/\allowbreak{}non-operational layer classification. & Record that the global guardrails are inherited. Phase-local no\_go\_guardrails may add stricter local constraints but may not replace, weaken, or scope away the global list. Any exception must be explicitly rejected unless it is already allowed by the global rule text and isolated outside the operational path. & No phase may be implemented or closed unless global\_no\_go\_inheritance is present in its scratchpad record and no local field contradicts the top-level global guardrails.\\
\hline
\addlinespace[0.24em]
Three-layer versioning policy & Any Lean source change;\newline any YAML planning/\allowbreak{}register change;\newline or any generator, importer, frontend, PDF/\allowbreak{}table rendering, navigation, API, or tool-code change. & metadata.lean\_data\_version + metadata.lean\_data\_timestamp;\newline metadata.planning\_data\_version + metadata.planning\_data\_timestamp;\newline metadata.tool\_infrastructure\_version + metadata.tool\_infrastructure\_timestamp;\newline change classification as Lean-code, planning/\allowbreak{}YAML state, or tool-infrastructure code;\newline latest-change overview row;\newline affected code-audit/\allowbreak{}guardrail/\allowbreak{}workflow rows. & Update exactly the affected version+timestamp pair(s), keep unaffected pairs unchanged, update compatibility aliases only as mirrors, and regenerate tables/\allowbreak{}TSV/\allowbreak{}PDF. Every row or export that records a semantic version must record the paired timestamp or be rejected as history-incomplete. & Latest-change overview shows all six semantic metadata fields;\newline generators validate all six;\newline SQL/\allowbreak{}Explorer snapshot contracts require all six;\newline no artifact claims a version without the matching timestamp.\\
\hline
\addlinespace[0.24em]
Pre-start feasibility and static-finding triage & Before starting any active phase after a user-reported green build, and before changing Lean code or closing an infrastructure gate. & active phase\_id;\newline current phase dependencies;\newline scratchpad preparation fields;\newline Section 2 static findings with Priority/\allowbreak{}Issue status/\allowbreak{}Addressing phase;\newline module manifest;\newline reuse/\allowbreak{}duplication audit;\newline object proof dossiers when relevant. & State whether the active phase is feasible as infrastructure-only, refactor, or Lean implementation. If a blocking issue is found, add or update a Static code finding row with R/\allowbreak{}Y/\allowbreak{}G priority, status, and addressing phase;\newline set or keep the phase red/\allowbreak{}yellow accordingly and backjump to the natural source layer rather than patching a consumer. & Implementation may start only after the active phase has a positive feasibility result or a precisely localized missing bridge. New findings may not remain prose-only;\newline they must be represented in Section 2 with a phase assignment and priority.\\
\hline
\addlinespace[0.24em]
Pre-Implementation Phase Check /\allowbreak{} semantic thought experiment & User asks to check the next semantically connected phases in advance, as a Gedankenexperiment, before implementation;\newline also applicable before a risky multi-phase refactor or derived-only bridge sequence. & Current active cursor;\newline phase table;\newline scratchpad;\newline object register and proof dossiers;\newline module implementation manifest;\newline reuse/\allowbreak{}duplication ledger;\newline consumption map;\newline static findings;\newline last user-reported green build when available;\newline requested semantic window or the next contiguous phase cluster inferred from dependencies. & For every inspected phase, create/\allowbreak{}update a pre-implementation check record and mark Y if the phase is implementable under current dependencies, or R if a dependency, object, proof route, module seam, source, or consumer-contract gap is missing. Add newly discovered problems to Static code findings. For every R record, generate one or more upstream addressing phases or a documented merge/\allowbreak{}backjump target before the blocked phase;\newline do not set the blocked phase active until those addressing phases are closed. & The check closes only when every phase in the semantic window has a Y/\allowbreak{}R result, each R result has a concrete addressing phase or documented merge target, generated phases have object-link policy and scratchpad stubs, and all changed records carry the current tool version/\allowbreak{}timestamp. Lean-code data timestamp remains unchanged unless Lean files or Lean inventory are changed.\\
\hline
\addlinespace[0.24em]
Concrete computable instance traversal & Before any operational Lean implementation starts, whenever a computable path, ExecComplex bridge, generated matrix/\allowbreak{}operator, recurrence, determinant, inverse candidate, or smoke-test route is introduced or modified. & active phase\_id;\newline target operational declarations;\newline at least one concrete finite instance or an explicit infrastructure-only justification;\newline expected executable witness;\newline current mathlib imports and known computability risks. & Fill concrete\_instance\_traversal, computable\_path\_not\_skeleton\_status, expected\_executable\_witness, mathlib\_complication\_log, noncomputable\_boundary\_audit, and mathlib\_mitigation\_or\_refactor\_layer in the scratchpad. Record whether the instance reaches the intended output object or where it stops. & Implementation may start only if the phase either has a concrete instance path through the computable code or is explicitly infrastructure-only. Any mathlib noncomputable leak must be isolated in an allowed mirror/\allowbreak{}non-operational layer or backjumped to the natural source layer;\newline it cannot silently enter the operational path.\\
\hline
\addlinespace[0.24em]
Phase preparation from reliable sources & User asks to prepare phase x with sources y,z or supplies established/\allowbreak{}reliable sources before implementation. & phase\_id;\newline source titles/\allowbreak{}URLs/\allowbreak{}files;\newline expected theory layer;\newline optional target objects. & Record source materials;\newline extract theory objects/\allowbreak{}formulas;\newline map them to CNNA object IDs;\newline update or add object rows;\newline split the phase if dependency depth is too high;\newline fill dependency audit, proof rehearsal, tactic plan, and natural fallback layer. Do not mark the phase green. & Phase remains Y/\allowbreak{}R/\allowbreak{}W according to dependency status. A prepared phase has a concrete positive proof path or a precisely localized missing bridge.\\
\hline
\addlinespace[0.24em]
Implementation from last green baseline & User asks to implement the active phase in Lean code. & last green package or last green CI snapshot;\newline active phase\_id;\newline prepared scratch-pad record;\newline object/\allowbreak{}phase dependencies;\newline applicable source notes;\newline explicit distinction between missing build evidence, verified red failure and verified green success. & Use positive construction first;\newline preserve import/\allowbreak{}consumer discipline;\newline add or modify only neutral code-surface modules;\newline record files changed, definitions, proof skeleton, tactics, forbidden-pattern checks and package result. If no build evidence is available, record artifact-generated /\allowbreak{} build required rather than red or green. & If the build is locally confirmed green or CEX15 receives verified green CI evidence, set phase G only through closeout evidence. If a local or CI build is verified red, set the attempted phase R and add/\allowbreak{}activate the missing repair bridge at the natural source layer. Absence of build evidence alone must not mark the phase red, green, closed or cursor-advanced.\\
\hline
\addlinespace[0.24em]
Object and quantities proof documentation workflow & A registered object or quantity is introduced, corrected from blue to public-final, exported through a handoff, or used as a proof-relevant code object. & object\_id;\newline current derivation site;\newline source objects;\newline CNNA definition;\newline intended proof target;\newline phase\_id;\newline semantic/\allowbreak{}API status. & Create or update the object proof dossier with source objects, definitions, proof goal, proof sketch/\allowbreak{}proof term reference, verification/\allowbreak{}audit result, documented yes/\allowbreak{}no, and link from the object register. Blue objects must record why proof-carrying is not yet semantically final. & An object may be public-final only when the object row and proof dossier agree: documented=yes, semantic status complete, API/\allowbreak{}handoff rights settled, and verification evidence points to the checked code or explicitly states that Lean build evidence is still local-only.\\
\hline
\addlinespace[0.24em]
Gap localization and backjump & During preparation or implementation a required object/\allowbreak{}proof bridge is missing. Also triggered automatically by red Lean CI through CEX15. & failed target statement;\newline currently available objects;\newline deepest successful dependency;\newline missing lemma/\allowbreak{}data shape. & Do not patch the consumer. Add or activate the earliest natural source phase where the structure should be generated;\newline record why the old phase cannot close;\newline preserve positive-path orientation. & The blocked phase is not closed. A smaller predecessor phase is created or activated with explicit source objects and proof goals.\\
\hline
\addlinespace[0.24em]
Post-green closeout & User reports the phase build is green or CEX15 receives a verified green CI build for the active phase. & compiler result;\newline exact zip/\allowbreak{}package used;\newline warnings or errors fixed;\newline current phase record;\newline for automated path: CI run ID, commit SHA, artifact manifest and scanner evidence;\newline documentation-output mode delta/\allowbreak{}full with selector or release checkpoint. & Fill implementation result;\newline update object statuses;\newline update linked phase-object map;\newline close scratch-pad record;\newline move active cursor to the next phase whose dependencies are complete. Automated closeout must use CEX15 transition invariants and may not advance the cursor on build success alone if evidence fields are incomplete. Regenerate/\allowbreak{}validate tables and emit the applicable documentation artifact: timestamp-bounded delta PDF for bounded planning changes, full PDF only for explicit full-release/\allowbreak{}full-audit checkpoints. & Exactly one active cursor remains. All green phases have complete implementation records;\newline generated tables and the requested documentation artifact validate/\allowbreak{}render. Automated green CI may not close a phase without CEX15 closeout evidence;\newline no workflow may require default full-PDF churn for bounded deltas.\\
\hline
\addlinespace[0.24em]
Local Lean build evidence rule & Any phase closeout or artifact package mentions build status. & Actual Lean build output from the current package or an explicit user report of a local green build. & Record whether the assistant generated artifacts only or whether a Lean build was actually run. Never infer green build from unchanged planning artifacts. & If no build was run, the package status is artifact-generated /\allowbreak{} local Lean build required, not green.\\
\hline
\addlinespace[0.24em]
Consumption classification and pruning discipline & A module/\allowbreak{}declaration appears unused, only BuildAll-reachable, status/\allowbreak{}audit-like, phase-semantic, or blue-object/\allowbreak{}API-unfinished. & module path;\newline declaration name/\allowbreak{}family;\newline import/\allowbreak{}reference evidence;\newline object ID if any;\newline proposed consumer or deletion target. & Classify the item in LEG5/\allowbreak{}LEG7/\allowbreak{}LEG8 before moving or deleting it. Decide keep/\allowbreak{}quarantine/\allowbreak{}wrap/\allowbreak{}reexport/\allowbreak{}archive/\allowbreak{}deletion-candidate. Update public API contract and firewall if it can be consumed by Phase 3+. & No deletion or public export until the item has a class, a consumer rule, and either a neutral API route or an explicit quarantine/\allowbreak{}deletion-candidate record. Deletion still requires later green build evidence.\\
\hline
\addlinespace[0.24em]
Pre-implementation module manifest & Before any Lean implementation phase starts or before an implementation prompt asks to add/\allowbreak{}modify code. & active phase\_id;\newline target objects;\newline current module inventory;\newline existing modules with matching semantics;\newline proposed new declarations;\newline import/\allowbreak{}firewall constraints. & Fill the scratchpad module manifest: exact count and names of proposed added modules and modified modules;\newline semantic placement rationale for each;\newline why existing modules are reused/\allowbreak{}extended or rejected;\newline public/\allowbreak{}legacy/\allowbreak{}API role;\newline expected import direction. & Implementation may start only if the module budget is explicit and compact. A new module is forbidden unless semantic fit, reuse, and import-firewall checks justify it.\\
\hline
\addlinespace[0.24em]
Reuse and duplicate-implementation audit & A planned declaration, helper, generator, ledger, request/\allowbreak{}output object, smoke test, or computation resembles existing code. & symbol/\allowbreak{}reference summary;\newline existing generator/\allowbreak{}API modules;\newline proposed target semantics;\newline object IDs and phase consumers. & Record whether the existing module is consumed, extended, wrapped, or rejected. If rejected, explain the semantic mismatch. Mark duplicate risk and update public API/\allowbreak{}manifest. & No duplicate implementation is allowed. The same mathematical/\allowbreak{}computational operation must have one canonical implementation path plus wrappers/\allowbreak{}reexports if needed.\\
\hline
\addlinespace[0.24em]
Canonical Smoke-Test Generator governance & A phase proposes generator-like code, experiment request/\allowbreak{}output code, smoke-test generation, or reusable theory-generation infrastructure. & existing Smoke-Test Generator modules and APIs;\newline target generator semantics;\newline phase/\allowbreak{}object consumers;\newline import constraints. & Route through the canonical generator seam whenever semantics match. If extension is needed, extend the canonical generator at its natural layer;\newline do not create a parallel generator. Register any genuine semantic split explicitly. & The canonical generator remains the single generator for the theory where applicable. Parallel generator modules are blocked unless the plan records a proven semantic split.\\
\hline
\addlinespace[0.24em]
Secondary long-term code-documentation and derivation-geometry workflow & A request proposes documenting the existing codebase, analyzing the geometry of the formal derivation graph, studying bottlenecks/\allowbreak{}dominators/\allowbreak{}source-sink paths, or reflecting on open-system emergence in the formalization topology. & Stable generated repo inventory;\newline object proof dossiers;\newline declaration consumption classification;\newline active not-before gates;\newline explicit derived-only/\allowbreak{}comparison-only boundary. & Register the goal as secondary if prerequisites are not closed. If active later, update coverage/\allowbreak{}dominators/\allowbreak{}layer/\allowbreak{}path diagnostics and link them to objects/\allowbreak{}modules without changing CNNA generative inputs. & Outputs are documentation, diagnostics, or comparison reports only. They may guide audits and refactors, but cannot become mathematical or physical premises without a separate CNNA-derived proof path.\\
\hline
\addlinespace[0.24em]
EXEC-2 rational boundary and mirror isolation & Any phase mentions ExecComplex, tau, sqrt(D), lattice embedding, Gram entries, eigenvalues, j-invariant, algebraic bridge, or mirror boundary. & Active phase;\newline target object IDs;\newline whether the value is rational/\allowbreak{}integer operational data or algebraic/\allowbreak{}analytic mirror data;\newline proof obligation for any proposed operational return value. & Fill exec\_complex\_boundary\_policy and noncomputable\_boundary\_audit. Update affected object ExecComplex/\allowbreak{}computable path text. Mark algebraic irrational data as mirror-only unless a rational/\allowbreak{}integer invariant is explicitly proved and typed. & No phase may introduce sqrt(D), tau, algebraic eigenvalues, or general algebraic complex coordinates into the operational Q x Q path. Operational return from mirror requires a stated invariant and proof target.\\
\hline
\addlinespace[0.24em]
Refactor regression and incremental migration & Any phase moves declarations, splits a module, introduces wrappers/\allowbreak{}aliases, quarantines legacy gates, or changes imports. & Module manifest;\newline declaration list or family;\newline current and target imports;\newline consumer list;\newline local build evidence plan. & Fill refactor\_regression\_protocol and refactor\_migration\_metrics. Record atomically delivered package boundary, declarations migrated count, import graph delta, consumer breakage audit, and local full-build status source. & A refactor phase closes only after the delivered package has local full-build evidence. Intermediate wrapper coexistence is allowed;\newline deleting legacy/\allowbreak{}status material before quarantine closure is not allowed.\\
\hline
\addlinespace[0.24em]
Lean/\allowbreak{}mathlib version management & Any phase proposes a new mathlib import, encounters noncomputable/\allowbreak{}classical pressure, or wants newer mathlib API than the frozen baseline. & Baseline version;\newline proposed import;\newline current compile/\allowbreak{}build evidence;\newline expected mathlib theorem/\allowbreak{}API;\newline operational vs mirror classification. & Fill lean\_mathlib\_version\_policy and mathlib\_complication\_log. Either stay on baseline, open a separate upgrade gate at a phase boundary, or re-route the dependency to mirror/\allowbreak{}non-operational code. & No silent version drift. No operational code may depend on a new mathlib API or noncomputable/\allowbreak{}classical pressure without an explicit gate and local build evidence.\\
\hline
\addlinespace[0.24em]
Phase-type-specific scratchpad validation & Every phase preparation or status change. & phase\_type;\newline phase status;\newline target code/\allowbreak{}object role;\newline active scratchpad. & Fill phase\_type and phase\_type\_specific\_contract. For refactor phases fill migration/\allowbreak{}regression fields;\newline for derivation phases fill proof/\allowbreak{}tactic/\allowbreak{}mathlib fields;\newline for comparison phases fill derived-only boundary;\newline for infrastructure phases justify absence of operational witness. & A phase cannot become active unless its phase\_type-specific evidence is prepared. Not-applicable filler is acceptable only with an explicit infrastructure-only or comparison-only justification.\\
\hline
\addlinespace[0.24em]
Governance-depth retrospective after first refactor & When Phase 1.1.2 or the first real code-change/\allowbreak{}refactor package closes. & Closed refactor report;\newline list of governance gates used;\newline mistakes prevented;\newline time/\allowbreak{}complexity overhead;\newline upcoming Phase 3+ structure. & Fill governance\_overhead\_review. Decide whether the current depth remains necessary or whether later phases should use a flatter structure with the same global guardrails. & Do not remove global no-go rules. Only the phase granularity may be flattened, and only after a recorded retrospective.\\
\hline
\addlinespace[0.24em]
Phase 3 to Phase 4 consumer-contract prebinding & Before preparing Phase 3.6 or Phase 4.1, or any wrapper around BoundaryOperatorProfile /\allowbreak{} ExecBoundaryMatrix. & BoundaryOperatorProfile shape;\newline Operational/\allowbreak{}ExecComplex consumer needs;\newline import direction;\newline generator policy. & Fill consumer\_contract\_direction. Phase 3.6 records the exported shape;\newline Phase 4.1 records how it consumes the shape without reconstructing Schur/\allowbreak{}DtN internals. & Phase 3 may not import Phase 4. Phase 4 may not duplicate generator/\allowbreak{}Schur/\allowbreak{}DtN logic;\newline it consumes the Phase-3 profile contract.\\
\hline
\addlinespace[0.24em]
Boundary-route extremal audit & Any phase that interprets IdealThread, CutSpec.full, Schur recursion, DtN boundary matrices, or string-like analogies. & Code references for IdealThread singleton carrier/\allowbreak{}family/\allowbreak{}coarsen lemmas;\newline CutSpec.full carrier theorem;\newline current generated boundary source;\newline planned consumer shape. & Record whether the phase consumes thread-scalar route, full-boundary route, or generated intermediate route. Mark analogical string language comparison-only unless a theorem gives a formal CNNA object. & No full matrix/\allowbreak{}operator, discriminant, spectrum, or physical interpretation may be promoted from the thread route unless the route and lifting theorem are explicitly present and code-checked.\\
\hline
\addlinespace[0.24em]
Horizontal/\allowbreak{}vertical Schur bridge workflow & A phase attempts to relate Cayley-Dickson/\allowbreak{}Hurwitz algebra to arithmetic discriminant, CM order, norm, or transfermatrix data through Schur complementation. & Closed S05/\allowbreak{}BR06 horizontal boundary route, closed O05/\allowbreak{}BR07 vertical arithmetic-shadow route, Transfermatrix objects TM0-TM2, CD7 Hurwitz boundary, and explicit module manifest for the bridge. & Create/\allowbreak{}update BR08/\allowbreak{}BR09/\allowbreak{}CD8 proof dossiers;\newline record commutation/\allowbreak{}naturality target, failure modes, concrete finite instance traversal, and no-input-backflow audit. & The bridge is either proved as a theorem over derived objects, split into smaller prerequisite lemmas, or marked blocked. It may not feed earlier phases unless a separate backreaction/\allowbreak{}input handoff explicitly permits it.\\
\hline
\addlinespace[0.24em]
Heuristic lattice test-matrix governance & Whenever Eisenstein/\allowbreak{}Gauss b-ary tree geometry, angle profiles, CM points, j-invariants, or L\_max sweeps are planned or implemented. & code-defined branching object/\allowbreak{}API;\newline L\_max set;\newline embedding kind;\newline angle profile;\newline EXEC-2 boundary status;\newline source objects L04/\allowbreak{}L06/\allowbreak{}O05/\allowbreak{}M05 when j/\allowbreak{}CM comparisons are used. & Record static/\allowbreak{}dynamic branch rows, concrete finite instance traversal, operational vs mirror fields, expected executable witness, and mathlib/\allowbreak{}noncomputable complications. Add or update L09/\allowbreak{}L10/\allowbreak{}M08/\allowbreak{}BR10 dossiers. & The test may close only if all rows consume the canonical generator/\allowbreak{}branching API, no duplicate tree generator is introduced, operational assertions are rational/\allowbreak{}integer/\allowbreak{}finite, and mirror-only CM/\allowbreak{}j comparisons cannot feed generator semantics.\\
\hline
\addlinespace[0.24em]
General early finite computation-test matrix & Before implementing or promoting any operational computable path, arithmetic wrapper, generator extension, lattice test, or refactor that claims computation is preserved. & Canonical generator/\allowbreak{}branching API;\newline finite instance parameters;\newline L\_max/\allowbreak{}branch-mode rows;\newline expected executable witness;\newline optional sentinel/\allowbreak{}comparison rows;\newline mathlib/\allowbreak{}noncomputable risk log. & Fill O07 or a consumer of O07;\newline update concrete\_instance\_traversal, expected\_executable\_witness, mathlib\_complication\_log, and no-go guardrails. Separate generic rows from sentinel rows. & The phase cannot close on hand-picked CM/\allowbreak{}Eisenstein/\allowbreak{}Gauss examples alone. At least one generic row family must traverse the actual computable path, or a documented failure/\allowbreak{}backjump is required.\\
\hline
\addlinespace[0.24em]
Regime-typed Heegner/\allowbreak{}CM finite test matrix workflow & Whenever finite computation tests, Eisenstein/\allowbreak{}Gauss sentinel rows, Heegner/\allowbreak{}class-number-one targets, CM/\allowbreak{}j targets, or noOuter/\allowbreak{}HasOuter regime claims are planned or implemented. & O07 generic finite rows;\newline derived regime split;\newline code-defined generator/\allowbreak{}branching API;\newline L\_max rows;\newline EXEC-2 boundary;\newline class-number data if Heegner interpretation is claimed. & Fill O08/\allowbreak{}N18/\allowbreak{}BR11/\allowbreak{}BR12 dossiers;\newline record regime tag, projection/\allowbreak{}complement-family status, generic-vs-sentinel status, operational witness, and mirror/\allowbreak{}comparison columns. Update mathlib/\allowbreak{}noncomputable log for any class-number/\allowbreak{}CM dependency. & The phase cannot close if h(D)=1 or a Heegner/\allowbreak{}j target is used to infer noOuter closure. Generic rows must exist;\newline noOuter-root sentinel rows and HasOuter/\allowbreak{}projection rows must be distinguished explicitly.\\
\hline
\addlinespace[0.24em]
Computable scalar architecture workflow & Any phase introduces scalar data beyond rational ExecComplex, uses mathlib Complex/\allowbreak{}Real/\allowbreak{}NumberField/\allowbreak{}QuadraticAlgebra, or proposes a Cayley-Dickson-derived C-layer. & O00/\allowbreak{}LEG20/\allowbreak{}LEG31, module manifest, reuse audit, expected executable witness, mathlib/\allowbreak{}noncomputable complication log, and CD noncircularity status if CD is involved. & Update LEG31/\allowbreak{}LEG32, O09/\allowbreak{}O10/\allowbreak{}BR13/\allowbreak{}BR14/\allowbreak{}CD9 dossiers;\newline classify each scalar as rational ExecComplex, exact quadratic schema, proof-carrying D-instantiation, integer/\allowbreak{}rational invariant, or analytic mirror-only. Record whether a new module is justified or existing module semantics suffice. & No new scalar layer may close without concrete finite instance traversal and proof-carrying projection back to allowed operational invariants. General Complex/\allowbreak{}Real and CD C-surfaces cannot become generator inputs, and exact quadratic instances require D-source provenance.\\
\hline
\addlinespace[0.24em]
Derived discriminant-indexed exact quadratic scalar workflow & Any phase creates or consumes exact quadratic coordinates, theta, sqrt(D)-like notation, CM tau coordinates, quadratic-field/\allowbreak{}order coordinates, or lattice rows requiring algebraic irrational data. & O09 schema, LEG32/\allowbreak{}GNG19, proof-carrying D-source from O05/\allowbreak{}O06 or equivalent ledger, module manifest/\allowbreak{}reuse audit, mathlib/\allowbreak{}noncomputable complication log. & Update O09/\allowbreak{}O10/\allowbreak{}BR14/\allowbreak{}BR13 dossiers;\newline record whether the phase is schema-only, D-instantiation, invariant projection, or mirror comparison. Record concrete finite D-ledger rows and allowed invariant exports. & No phase may close with a free sqrt(D), an axiom, or an analytic Complex/\allowbreak{}CD source for theta. If D is unavailable or semantically blue, backjump to the discriminant-source correction rather than patching the scalar consumer.\\
\hline
\addlinespace[0.24em]
Phase-path dependency and object placement audit & Before any implementation phase after Phase 4, or when moving/\allowbreak{}wrapping an object from N/\allowbreak{}L/\allowbreak{}M/\allowbreak{}SP/\allowbreak{}P/\allowbreak{}TM/\allowbreak{}BR/\allowbreak{}CD groups. & Object register, source-object list, target module manifest, dependency-depth audit, and current phase-path classification. & Update phase title, object Phase field, proof dossier, scratchpad dependency\_depth\_audit, consumer\_contract\_direction, and no-go notes. Record whether the object belongs to Path A, Path B, or a late explicit bridge. & Phase may close only if its title, object group and dependencies agree. Finite spectral/\allowbreak{}physical consumers cannot be blocked behind modular/\allowbreak{}j/\allowbreak{}analytic phases unless a proof-carrying bridge says so.\\
\hline
\addlinespace[0.24em]
Phase-object-link policy workflow & Whenever a phase is added, split, merged, moved, or whenever an object Phase field changes. & Phase hierarchy;\newline object register;\newline direct object ownership map;\newline child-phase scope;\newline moved/\allowbreak{}audit/\allowbreak{}closeout status. & Update phase Object link policy, Object link explanation, Child object scope, object Phase fields, proof dossiers and phase table generator if needed. & The generated phase table may not contain an unexplained empty Linked objects cell. Direct ownership is never inferred from child scope;\newline child-scope summaries are display-only.\\
\hline
\addlinespace[0.24em]
Ontic subsystem-minimality workflow & Whenever a phase introduces or consumes dynamics, spectra, operator algebra, geometry, arithmetic, stochasticity, open-system language, unitarity, AQFT/\allowbreak{}physics interpretation, or environment/\allowbreak{}complement effects. & POST1/\allowbreak{}LEG35/\allowbreak{}GNG22;\newline source-object chain from IDEAL\_infty formal background to subsystem/\allowbreak{}cut/\allowbreak{}approximant, boundary, regime, complement families and Schur/\allowbreak{}DtN or later certified bridge. & Update object proof dossier, scratchpad source\_to\_cnna\_mapping, dependency\_depth\_audit, no-go guardrails, and consumer-contract direction. Mark comparison-only interpretations explicitly and prohibit them from supplying definitions. & A phase may close only if non-subsystem structures are shown as derived consequences or comparison-only interpretations. Any primitive import of dynamics/\allowbreak{}physics/\allowbreak{}arithmetic/\allowbreak{}geometry/\allowbreak{}environment forces a backjump to the natural subsystem/\allowbreak{}complement generating layer.\\
\hline
\addlinespace[0.24em]
Lean-focus implementation-packet workflow & Before a Lean implementation phase starts, or when the user asks for the next semantically connected phases to be checked with minimal cognitive/\allowbreak{}tooling overhead. & Active phase, next semantic phase window, object register, scratchpad, module manifest, static findings, import/\allowbreak{}declaration inventory, and current pre-implementation-check ledger. & Generate a compact implementation packet and avoid full PDF/\allowbreak{}declaration dumps unless the phase-check or audit-full mode explicitly requires them. Record G/\allowbreak{}Y/\allowbreak{}O/\allowbreak{}R thought-experiment results, predecessor assumptions, proof risk, and any red addressing phase proposals. & Lean coding may start only from a compact packet that names relevant modules, producer/\allowbreak{}consumer objects, red blockers, orange predecessor assumptions, yellow active implementability, forbidden shortcuts, and comparison-only theory boundaries.\\
\hline
\addlinespace[0.24em]
Local theory context versus end-theory comparison workflow & Whenever a phase or module mentions established mathematical/\allowbreak{}physical theory, mathlib background, or a large target theory such as AQFT/\allowbreak{}OQS/\allowbreak{}modular forms/\allowbreak{}CM/\allowbreak{}Moonshine. & Module imports, declared CNNA objects, object proof dossiers, local background-theory context, and end-theory comparison ledger. & Classify local established theory as formalization/\allowbreak{}proof-method/\allowbreak{}library context;\newline classify large end theories only as comparison/\allowbreak{}target/\allowbreak{}consistency horizons. Record that neither role is an ontic CNNA generator unless derived internally. & A module may use established mathematics as formal language, but a CNNA object may close only if its producer is internal provenance. Large end theories remain comparison-only until CNNA derives a bridge object.\\
\hline
\addlinespace[0.24em]
DB-REVIEW database-ready documentation and expert-review workflow & Whenever the YAML schema, generated views, phase/\allowbreak{}object records, proof dossiers, theory/\allowbreak{}reference markings, online/\allowbreak{}interactive documentation, or database/\allowbreak{}export plans are edited. & Lean-code data timestamp;\newline YAML master;\newline phase/\allowbreak{}object/\allowbreak{}proof-dossier records;\newline status legend;\newline global no-go guardrails GNG23-GNG24;\newline local theory context/\allowbreak{}end-theory comparison ledgers when applicable. & Keep Lean as primary truth system, generated views as secondary, and expert review as mandatory. Record whether a row is database-ready-later, review-required, blue-audit-needed, lean-primary, or generated-secondary when those semantics matter. Preserve round-trip-friendly IDs and explicit phase/\allowbreak{}object/\allowbreak{}review links. & No generated documentation may close by declaring semantic infallibility or 100-percent truth. Missing blue/\allowbreak{}reference markings must remain visible as audit debt, and database/\allowbreak{}export work must not become a CNNA-generative input.\\
\hline
\addlinespace[0.24em]
Context-filtered documentation slice generation & User asks for documentation about a specific object/\allowbreak{}topic, or a phase closeout requires documenting only the changed context rather than regenerating the whole PDF for human review. & Natural-language query or curated profile id;\newline current Single-YAML-Master;\newline regenerated generated\_exports/\allowbreak{}*.tsv;\newline optional slug. & Run scripts/\allowbreak{}generate\_tables.py first;\newline run scripts/\allowbreak{}generate\_context\_doc.py with query/\allowbreak{}profile;\newline package context\_docs/\allowbreak{}<slug>/\allowbreak{} manifest, section summary, TeX/\allowbreak{}PDF, and compile log. Add/\allowbreak{}adjust curated profile terms in YAML if the topic is recurring. & Every major documentation section is present in the context packet;\newline each section records matched/\allowbreak{}total row counts;\newline output states generated-secondary status, Lean-primary rule, and expert-review requirement.\\
\hline
\addlinespace[0.24em]
Semantic module consolidation and compatibility-preserving refactor workflow & Before new substantive Lean code after governance closure, or whenever the user asks whether semantically related modules should be merged, compacted, facaded, or kept separate. & LEG10 module manifest;\newline LEG11 reuse ledger;\newline LEG5 consumption map;\newline repo\_inventory/\allowbreak{}raw\_export modules/\allowbreak{}import\_edges/\allowbreak{}symbol references;\newline candidate file list;\newline direct code spot-checks;\newline current public/\allowbreak{}API and proof-dossier records. & Update LEG36, code\_audit findings, reuse\_duplication\_ledger, scratchpad refactor fields, proposed modules, duplicate-risk symbols, natural-layer placement, facade/\allowbreak{}compatibility plan, and regression metrics. Set downstream content phases to orange if this gate must precede them. & No Lean merge may close without old-module facade or explicit migration plan, import graph delta, declaration/\allowbreak{}consumer delta, local full-build evidence, and proof that semantic/\allowbreak{}provenance auditability is preserved or improved.\\
\hline
\addlinespace[0.24em]
Root artifact retention and generated-output cleanup workflow & Before packaging a tool release, after full/\allowbreak{}context documentation generation, or when root PDFs/\allowbreak{}TeX/\allowbreak{}Markdown reports accumulate. & Current Repository root file list;\newline Single-YAML-Master;\newline README.md;\newline generated tables/\allowbreak{}exports;\newline optional generated\_docs/\allowbreak{}full output;\newline archive/\allowbreak{}reports policy. & Keep README.md in root;\newline move historical Markdown reports to archive/\allowbreak{}reports/\allowbreak{};\newline delete disposable root PDF/\allowbreak{}TeX/\allowbreak{}LaTeX auxiliary files;\newline ensure full documentation is generated under generated\_docs/\allowbreak{}full/\allowbreak{} on demand;\newline update DBR8/\allowbreak{}GNG27 if classification changes. & Project root contains no retained generated PDF/\allowbreak{}TeX/\allowbreak{}LaTeX build products and no active historical Markdown reports beyond README.md. All retained source data remains in YAML/\allowbreak{}scripts/\allowbreak{}tables/\allowbreak{}generated\_exports/\allowbreak{}repo\_inventory/\allowbreak{}repo\_snapshot as appropriate.\\
\hline
\addlinespace[0.24em]
Root artifact retention and Lean/\allowbreak{}Lake build-entry preservation workflow & Before packaging a tool release, after root cleanup, or when Lean/\allowbreak{}Lake build-control files are added, removed, or moved. & Package root file list;\newline CNNA.lean;\newline lakefile.lean;\newline lake-manifest.json;\newline lean-toolchain;\newline Repository/\allowbreak{}repo\_snapshot/\allowbreak{}CNNA source tree;\newline README.md;\newline cleanup script policy. & Ensure package-root Lean/\allowbreak{}Lake project-control files are preserved;\newline ensure generated docs/\allowbreak{}reports remain outside active root;\newline update DBR10/\allowbreak{}GNG28 and README if the build-entry surface changes. & Package root contains the source/\allowbreak{}control files needed for local Lake build and contains no retained generated PDF/\allowbreak{}TeX/\allowbreak{}LaTeX byproducts or historical active reports.\\
\hline
\addlinespace[0.24em]
Root license, README and credit-surface workflow & Before packaging a tool release, when licensing/\allowbreak{}credits change, or when README files appear outside the package root. & Package root file list;\newline README.md;\newline LICENSE;\newline license\_notice table;\newline metadata editor/\allowbreak{}assistant/\allowbreak{}minimum\_assistant\_requirement;\newline full/\allowbreak{}context documentation generators;\newline cleanup policy. & Keep README.md and LICENSE in package root;\newline remove duplicate active README files in subdirectories;\newline update license\_notice/\allowbreak{}Stats/\allowbreak{}PDF author metadata;\newline update DBR11/\allowbreak{}GNG29 and README if credits, license scope, or tested assistant setup changes. & Exactly one maintained README.md exists at package root;\newline root LICENSE exists;\newline generated full/\allowbreak{}context PDFs expose Editor and Assistant credits;\newline package contains no root PDF/\allowbreak{}TeX/\allowbreak{}LaTeX byproducts.\\
\hline
\addlinespace[0.24em]
Generated PDF section-pagebreak workflow & When full or context documentation layout requirements change, especially section-boundary/\allowbreak{}page-navigation rules for review PDFs. & Full-document generator, context-document generator, Master YAML metadata, representative full/\allowbreak{}context PDF outputs, and the user-facing README command surface. & Patch generator preambles or layout macros;\newline keep disposable generated TeX/\allowbreak{}PDF out of the package root;\newline regenerate tables and representative PDFs;\newline update DBR12/\allowbreak{}GNG30/\allowbreak{}SCF49 and README if user commands or layout expectations change. & The full documentation PDF and at least one context documentation PDF compile successfully, their generated TeX contains the section-pagebreak hook, and no Lean source or context-matching semantics are changed.\\
\hline
\addlinespace[0.24em]
Phase status release and red-lock workflow & Before adding phases, changing a phase status, advancing the active cursor, or preparing a Lean-derivation implementation prompt. & phase row;\newline implementation\_scratchpad.phase\_type;\newline prerequisite object IDs;\newline proof dossier status;\newline pre-implementation check result;\newline current active cursor;\newline relevant static findings and guardrails. & If phase\_type=derivation and the phase is not closed, active, or explicitly implementable-after-predecessors, set it red. To release red, add a pre-implementation check that references already-derived Lean objects and sets yellow only for the unique active cursor or orange for later implementable phases. Keep white only for non-direct-Lean refactor/\allowbreak{}tool/\allowbreak{}infrastructure/\allowbreak{}documentation/\allowbreak{}comparison phases. & Generated tables pass validation: exactly one yellow phase;\newline zero white derivation phases;\newline red derivation phases carry lock language;\newline any yellow/\allowbreak{}orange release has explicit prerequisite evidence or a corresponding pre-implementation check record.\\
\hline
\addlinespace[0.24em]
Documentation section explanation and appendix workflow & When a documentation section/\allowbreak{}table is confusing, redundant, or appears to duplicate a primary control surface. & Full-document section order, context-document SECTION\_ORDER, static table paragraphs, Master YAML sections, generated exports, and user review feedback. & Add or tighten static generator-owned explanations, keep table rows dynamic, move narrow technical tables to appendix when they are not primary reading surfaces, and update DBR/\allowbreak{}GNG/\allowbreak{}static finding records if semantics change. & A reader can identify what each section/\allowbreak{}table displays, where the data comes from, why it is shown, and which table is authoritative for phase release. No dynamic prose rows are added only to explain static section purpose.\\
\hline
\addlinespace[0.24em]
Pre-substantive-code readiness and proof-gap audit workflow & Before starting any new substantive Lean/\allowbreak{}content implementation, especially before activating 1.1.1.2 or any later red locked derivation phase. & Current full zip;\newline repo\_snapshot/\allowbreak{}CNNA source tree;\newline repo\_inventory/\allowbreak{}raw\_export module/\allowbreak{}import/\allowbreak{}declaration inventory;\newline generated\_exports consumption/\allowbreak{}object/\allowbreak{}proof/\allowbreak{}phase tables;\newline local build result if refactor is executed. & Fill LEG37 source-intake ledger;\newline create LEG38 neutral naming/\allowbreak{}facade plan;\newline execute only approved LEG39 refactor package;\newline regenerate inventory and close LEG40;\newline issue LEG41 release proof or create upstream red blockers. & Exactly one active Y phase remains. No new substantive Lean code begins until LEG41 references closed derived-only/\allowbreak{}verified objects and the stable neutral API/\allowbreak{}source surface.\\
\hline
\addlinespace[0.24em]
Provider-neutral SQL public read-only mirror and web-frontend workflow & When the database/\allowbreak{}web-publication sidechain is activated or when a provider-neutral SQL database mirror, public read-only API, cnna-project.net frontend, reviewer annotation layer or direct importer is designed. & Master YAML, generated TSV exports, repo\_inventory/\allowbreak{}raw\_export, all three timestamp-complete version pairs (lean\_data\_version/\allowbreak{}lean\_data\_timestamp, planning\_data\_version/\allowbreak{}planning\_data\_timestamp, tool\_infrastructure\_version/\allowbreak{}tool\_infrastructure\_timestamp), source zip name, YAML/\allowbreak{}repo inventory/\allowbreak{}tool hashes, public read-only database endpoint/\allowbreak{}credential for read-only frontend, local admin write credential key only via .env for admin importer if/\allowbreak{}when executed, and CEX16 red-failure-history schema when CI history is exposed. lean\_data\_version is the Lean toolchain/\allowbreak{}API level;\newline lean\_data\_timestamp is the Lean source/\allowbreak{}data snapshot time. & Add or update DBR15-DBR20, schema/\allowbreak{}import/\allowbreak{}read-only access policy/\allowbreak{}frontend contracts, snapshot/\allowbreak{}checksum records, read-only grant/\allowbreak{}access tests, public view definitions, README/\allowbreak{}manual commands, archive report and failure-history views if red CI history is public. Never store secrets in repo/\allowbreak{}chat. & A DB mirror phase may close only when every public current row traces to a snapshot\_id with all three version+timestamp pairs, public read-only access is read-only through api.* views, internal/\allowbreak{}admin schemas are closed, admin write credential remains server/\allowbreak{}local-admin only, and a YAML/\allowbreak{}export roundtrip or import report is archived. If CI history is exposed, red failure rows must be visibly separated from green current rows and linked to repair phases. The Lean pair must remain asymmetric when appropriate: toolchain version is not the source snapshot timestamp.\\
\hline
\addlinespace[0.24em]
Generated PDF typography profile workflow & When generated full/\allowbreak{}context documentation readability or typography scaling changes are requested. & Full-document generator, context-document generator, table generator, representative full/\allowbreak{}context PDF outputs, generated TeX audit and README command surface. & Patch generator-level size profile, regenerate tables, compile representative PDFs, render sample pages, update DBR21/\allowbreak{}GNG35/\allowbreak{}SCF54 and README if user-facing layout expectations change. & Generated full and context PDFs compile with the intended typography profile, generated TeX shows the larger size profile, representative pages render without obvious clipping/\allowbreak{}black-glyph failure, and no Lean/\allowbreak{}YAML semantics change.\\
\hline
\addlinespace[0.24em]
Context-class assignment & Whenever a phase or object is created, moved, reclassified, exported to DB, or rendered in a phase/\allowbreak{}object table. & Declared Context class from context\_class\_taxonomy plus phase\_type/\allowbreak{}object provenance/\allowbreak{}status fields. & Master YAML phase/\allowbreak{}object row;\newline proof dossier if object;\newline scratchpad if phase;\newline generated tables and TSV exports. & No blank Context class;\newline class must not contradict handoff/\allowbreak{}access rights or generated table semantics.\\
\hline
\addlinespace[0.24em]
Object patch and supersession lifecycle workflow & A Lean/\allowbreak{}planning/\allowbreak{}tool object is patched, semantically replaced, or receives a corrected canonical description/\allowbreak{}dossier. & Old object ID, old dossier pointer, new or corrected object ID/\allowbreak{}dossier, patching phase ID, reason for replacement, snapshot/\allowbreak{}version timestamps and affected consumers. & Do not overwrite the old object row. Set old Object lifecycle state to superseded, fill Lifecycle successor/\allowbreak{}current object and Lifecycle note, keep old proof dossier readable, create or update the successor object/\allowbreak{}dossier, link the successor phase to the old phase/\allowbreak{}object, and record consumer migration duties. & The lifecycle chain is acyclic, the old object remains queryable as historical provenance, the new/\allowbreak{}current object has a current dossier pointer, and no phase is marked outdated.\\
\hline
\addlinespace[0.24em]
Object rename lifecycle workflow & A Lean declaration/\allowbreak{}module/\allowbreak{}object or planning object is renamed without intending to preserve the old name as current canonical API. & Old name/\allowbreak{}ID, new name/\allowbreak{}ID, rename phase, module/\allowbreak{}declaration diff evidence or planning diff, affected imports/\allowbreak{}consumers, old and new dossier pointers. & Retain the old object row as renamed-alias, point it to the new/\allowbreak{}current object, preserve the old dossier as historical alias evidence, create/\allowbreak{}update the new object row and dossier, and record any migration links in module/\allowbreak{}object registers. & Explorer/\allowbreak{}SQL views can resolve old aliases to the current object while still displaying original provenance;\newline no historical row is deleted or silently merged.\\
\hline
\addlinespace[0.24em]
Object deletion or retirement lifecycle workflow & An object/\allowbreak{}declaration/\allowbreak{}module is removed from Lean or planning, or intentionally retired as no longer needed. & Removed object ID, deletion/\allowbreak{}retirement phase, deletion evidence, reason, downstream consumers, replacement object if any, snapshot/\allowbreak{}version timestamps. & Retain the object row with Object lifecycle state deleted-retired, keep its dossier, fill successor if a replacement exists or explain terminal retirement in Lifecycle note, and create migration/\allowbreak{}blocker tasks for affected consumers. & The object no longer appears as current in Explorer projections but remains visible in history/\allowbreak{}provenance queries;\newline successor chains terminate in a current object or explicit terminal retirement.\\
\hline
\addlinespace[0.24em]
Phase patch-link workflow & A later correction/\allowbreak{}refactor/\allowbreak{}patch changes or replaces an object that was produced or consumed by an earlier phase. & Old phase ID, new patch/\allowbreak{}successor phase ID, affected objects, reason for patch, version/\allowbreak{}timestamp evidence and migration obligations. & Never mark the old phase outdated. Keep it as a historical implementation event;\newline link the new phase to the old phase via patch/\allowbreak{}successor relation in phase notes, object lifecycle rows and scratchpad/\allowbreak{}dossier references. & Both phases remain independently queryable;\newline the new phase explains the correction and the old phase remains part of historical provenance without being treated as current object validity.\\
\hline
\addlinespace[0.24em]
Dossier and scratchpad ownership workflow & A question arises whether evidence belongs under a phase scratchpad, object dossier, lifecycle ledger or generated history entry. & Entity type, evidence type, target object/\allowbreak{}phase IDs, version/\allowbreak{}timestamp and whether the evidence describes implementation execution, object validity, lifecycle migration or generated diff history. & Store implementation process evidence in the phase scratchpad;\newline store object validity/\allowbreak{}proof evidence in the object dossier;\newline store patch/\allowbreak{}rename/\allowbreak{}delete/\allowbreak{}supersession evidence in object lifecycle fields and lifecycle ledger;\newline store snapshot diffs in history records. Cross-link rather than duplicate or overwrite. & Generated views can reconstruct both the historical implementation path and the current object state without contradiction.\\
\hline
\addlinespace[0.24em]
Lean declaration total-classification workflow & Phase 21.0, LEG37, Explorer schema work or any public API/\allowbreak{}refactor phase needs to claim that the current Lean declaration surface is fully understood. & Current Lean snapshot timestamp, lean toolchain version, repo\_inventory/\allowbreak{}raw\_export, declaration\_consumption\_records.generated.tsv, object register, proof dossiers, lifecycle ledger and direct source scan evidence. & For every Lean declaration, record exactly one declaration\_kind, semantic\_role, triviality\_tier, significance/\allowbreak{}documentation policy, lifecycle\_state, public\_surface\_status and owning CNNA object or explicit non-object bucket. Unknown/\allowbreak{}manual-review values may be stored only as open blockers with natural-origin follow-up action. & LEG37/\allowbreak{}CEX0/\allowbreak{}21.0 may close only when all declarations have total classification or every unresolved declaration is explicitly routed to a red/\allowbreak{}orange upstream phase;\newline consumption class alone never suffices.\\
\hline
\addlinespace[0.24em]
Lean declaration triviality override workflow & An automatic scanner marks a declaration trivial/\allowbreak{}nontrivial based on rfl, projection shape, theorem/\allowbreak{}structure keyword, module path or consumer count. & Automatic candidate label, proof/\allowbreak{}body/\allowbreak{}consumer evidence, source object or owner bucket, handoff/\allowbreak{}boundary/\allowbreak{}API relevance and manual reviewer note when the automatic label is overridden. & Store both the automatic candidate and the curated final triviality/\allowbreak{}significance/\allowbreak{}documentation decision. In particular, rfl/\allowbreak{}projection declarations that encode handoff, boundary or fundamental definitional contracts may be promoted;\newline ledger/\allowbreak{}status structures may be demoted when they carry no mathematical payload. & Explorer and proofbook views render the curated final class plus override reason when present;\newline no heuristic-only triviality class may release a declaration to public-final or exclude it from required documentation.\\
\hline
\addlinespace[0.24em]
Lean declaration owner-bucket and lifecycle workflow & A declaration is added, patched, renamed, deleted, retired, superseded, or cannot be attached to a registered object. & Declaration name, module path, previous/\allowbreak{}current snapshot IDs, object/\allowbreak{}dossier candidate, successor/\allowbreak{}predecessor relation if any, reason for non-object bucket if not owned by a CNNA object, and public-surface status. & Attach the declaration to exactly one owning object or explicit non-object bucket, assign lifecycle state, preserve old names as historical rows when renamed/\allowbreak{}deleted/\allowbreak{}retired, and update successor/\allowbreak{}current pointers without overwriting old evidence. & Generated history can reconstruct both current declaration state and original provenance;\newline lifecycle/\allowbreak{}successor chains remain acyclic and no active consumer points to a deleted-retired or superseded declaration except through an explicit historical/\allowbreak{}facade edge.\\
\hline
\addlinespace[0.24em]
Original roadmap integration workflow & An uploaded legacy/\allowbreak{}original roadmap contains content not yet represented in the Single-YAML-Master planning state. & Roadmap source file, current master YAML, object register, phase table, scratchpads, dossiers, external-source cross-check, guardrails and lifecycle records. & Copy the source into archive/\allowbreak{}source\_materials when useful;\newline add only missing source hierarchy, gate, chain, migration, de-seeding, register, context or future-target obligations. Route accepted items to explicit phases/\allowbreak{}objects/\allowbreak{}dossiers/\allowbreak{}cross-check rows;\newline mark rejected/\allowbreak{}deferred items by reason. Do not import roadmap claims as axioms or checked Lean facts. & Every accepted roadmap obligation has a phase/\allowbreak{}object/\allowbreak{}dossier/\allowbreak{}cross-check row or an explicit no-action reason;\newline Lean code remains unchanged unless a separate implementation phase is executed.\\
\hline
\addlinespace[0.24em]
De-seeding and old-path regression workflow & A Seed, Scaffold, Window, Boundary-first or legacy module is proposed for current canonical API or de-seeded naming. & Source map, derived-only gate evidence, notation/\allowbreak{}readability binding, free-parameter audit, old-path regression oracle, importgraph cycle check, rollback snapshot and affected object lifecycle rows. & Create/\allowbreak{}close the natural source-map gate before renaming;\newline keep old object/\allowbreak{}module as regression or lifecycle-retired provenance;\newline update object lifecycle/\allowbreak{}dossiers and module manifest;\newline do not patch the consumer phase. & The new current object has source-map and gate evidence;\newline the old path remains queryable;\newline importgraph is acyclic;\newline rollback snapshot exists;\newline no free legacy parameter remains hidden.\\
\hline
\addlinespace[0.24em]
Roadmap chain-sheet to phase/\allowbreak{}object workflow & A roadmap chain sheet identifies a multi-step dependency path such as modular B core, ParameterClosure or D/\allowbreak{}E bridge. & Chain definitions, prerequisites, proof/\allowbreak{}claim text, target modules, current plan coverage and missing source objects. & Represent the chain as a phase/\allowbreak{}object/\allowbreak{}gate group with explicit predecessor order;\newline add guardrails where shortcuts are forbidden;\newline mark as reference-only until CNNA source objects close. & The chain is visible in phase order, object register, cross-check rows and scratchpad obligations without implying that the chain is already proved.\\
\hline
\addlinespace[0.24em]
REALOQS legacy code scan workflow & A legacy REALOQS source archive is uploaded or re-scanned. & versioned source archive;\newline scan timestamp;\newline current planning package;\newline target focus such as AQFT/\allowbreak{}PillarB & Archive source, generate module/\allowbreak{}declaration metrics, record noncomputable/\allowbreak{}simp/\allowbreak{}sorry scan, add external-source rows, update RLC ledgers and lifecycle pointers. & A scan report exists and every accepted finding is routed to reference/\allowbreak{}migration ledgers, not to CNNA truth.\\
\hline
\addlinespace[0.24em]
REALOQS AQFT migration-disposition workflow & A legacy AQFT declaration/\allowbreak{}module is considered for CNNA reuse. & legacy module path;\newline declaration list;\newline dependency/\allowbreak{}import context;\newline current CNNA source objects and handoff contract & Classify as target-profile, semantic-rewrite, mathlib replacement, regression-only, direct-port-forbidden, reject/\allowbreak{}defer;\newline update object dossier and phase gate. & No implementation starts until each accepted target has a CNNA source-map, neutral target module and forbidden-dependency audit.\\
\hline
\addlinespace[0.24em]
REALOQS declaration crosswalk workflow & Legacy declarations need Explorer visibility or migration analysis. & RLC0 scan data, CEX11 taxonomy, current CNNA declaration inventory, lifecycle ledger & Map each legacy declaration to owner bucket, semantic role, target CNNA object/\allowbreak{}phase or reject/\allowbreak{}defer bucket;\newline mark lifecycle and edge type as legacy/\allowbreak{}reference. & Explorer may show the legacy node only if its non-CNNA epistemic status and edge semantics are explicit.\\
\hline
\addlinespace[0.24em]
Legacy regression-harness extraction workflow & A REALOQS example/\allowbreak{}meta proof is useful as a future test. & legacy example/\allowbreak{}meta module, target CNNA phase, expected old behavior, source-map risk & Create regression expectation without importing code;\newline quarantine example;\newline link to future derived replacement and rollback/\allowbreak{}compatibility check. & Regression harness is recorded as test evidence only, not as a source object or closure proof.\\
\hline
\addlinespace[0.24em]
Generated-data derived-only closure workflow & A phase/\allowbreak{}object claims derived-only closure, arithmetic/\allowbreak{}physics seed success, or green status for a nontrivial path. & Structural skeleton;\newline local proof obligations;\newline generated data source;\newline operational smoke path;\newline verification evidence;\newline current Lean data timestamp. & Record skeleton, proof modules, generated-data origin, smoke path, and verification result in scratchpad/\allowbreak{}dossier;\newline mark claim yellow/\allowbreak{}orange/\allowbreak{}red until the complete path is shown. & The phase/\allowbreak{}object can close only when the operational generated-data path is explicit and the dossier distinguishes skeleton success from real data validation.\\
\hline
\addlinespace[0.24em]
Natural-origin missing-structure workflow & Implementation, proof search or pre-implementation thought experiment discovers a missing structure in a downstream phase. Includes red CI failures produced by the active phase branch/\allowbreak{}PR. & Failing consumer, missing structure description, candidate source layers, rejected patch sites, dependency graph, proof/\allowbreak{}audit reason for the selected origin. & Mark downstream phase red;\newline add static finding;\newline create/\allowbreak{}activate the earliest natural upstream phase;\newline document why the structure belongs there and not at the consumer. For red CI, also expose the failure through the visible failure-history channel and link it to the generated repair/\allowbreak{}localization phase. & Downstream phase can reopen only after the upstream natural-origin phase closes and lifecycle/\allowbreak{}provenance links are updated.\\
\hline
\addlinespace[0.24em]
Implementation-thought-experiment phase-readiness workflow & The user asks whether future phases are implementable or which positive path phase should become active. & Current phase/\allowbreak{}object table, dependencies, linked objects, module seams, no-go guardrails, implementation packet, known red blockers. & Simulate implementation pressure rather than only checking listed dependencies;\newline record hidden dependencies, route blockers upstream, and assign G/\allowbreak{}Y/\allowbreak{}O/\allowbreak{}R status with evidence. & A phase is yellow only when the thought experiment finds no missing source object/\allowbreak{}consumer contract;\newline orange only with explicit predecessor closures;\newline red if an upstream natural-origin blocker exists.\\
\hline
\addlinespace[0.24em]
Explorer/\allowbreak{}YAML epistemic-control workflow & CNNA planning complexity exceeds linear-document manageability or a new visual/\allowbreak{}database presentation is proposed. & Master YAML, generated exports, stable IDs, typed nodes/\allowbreak{}edges, lifecycle status, declaration classification, external-reference/\allowbreak{}hypothesis markers. & Treat YAML/\allowbreak{}SQL/\allowbreak{}Explorer as epistemic control and review surfaces;\newline keep PDF/\allowbreak{}TeX as generated renderings;\newline ensure all Explorer status displays are backed by machine-readable fields. & Explorer/\allowbreak{}UI may present the project only when it preserves Lean/\allowbreak{}YAML truth separation, edge semantics, lifecycle status, and review/\allowbreak{}hypothesis markers.\\
\hline
\addlinespace[0.24em]
Pillar co-development and acyclic handoff workflow & A later plan requires B, C, D and E to be developed in a mutually constraining way from Pillar A. & Pillar-specific source objects, A->X handoffs, input handoff rights, backreaction channels, graph edge types, cycle checks. & Represent co-development as explicit acyclic handoff/\allowbreak{}backreaction windows;\newline do not replace the partial order by informal mutual dependency. & The plan may show co-development visually only when proof-producing provenance remains acyclic and all backreaction enters through declared input handoffs.\\
\hline
\addlinespace[0.24em]
GitHub backup-gated assistant/\allowbreak{}repository automation workflow & The assistant or an automated script is asked to modify the CNNA GitHub repository rather than only an uploaded zip. & Repository full name, current default-branch SHA, branch-protection rules, GitHub App/\allowbreak{}API permissions, intended phase/\allowbreak{}object scope and active no-go guardrails. & Create or verify a backup tag/\allowbreak{}branch for the current default-branch SHA before any mutable work;\newline create a work branch;\newline commit changes there;\newline open a PR;\newline run required CI;\newline record backup ref, work branch, PR, commit SHA and CI status in the planning/\allowbreak{}export evidence. If backup ref creation or verification fails, stop before work-branch mutation. & Normal path exits only through PR plus green required checks. Direct main writes are blocked except explicit emergency override with recorded backup, reason and rollback path. No automated repository write may proceed without a recoverable backup ref for the parent commit it modifies.\\
\hline
\addlinespace[0.24em]
CI snapshot-artifact and Webhosting-MySQL import workflow & A green Lean/\allowbreak{}inventory build should update the public Explorer or MySQL-backed production/\allowbreak{}current mirror;\newline red builds must instead route to CEX15/\allowbreak{}CEX16 and publish visible failure history without current swap. & Pinned lean-toolchain/\allowbreak{}lake-manifest, Master YAML, generated exports, repo inventory, snapshot\_manifest schema, Webhosting import endpoint or server-local importer, DB credentials held outside public code. & Run lake build and scans;\newline generate YAML/\allowbreak{}JSON/\allowbreak{}TSV/\allowbreak{}SQL artifacts plus hashes. On green CI, upload artifact, import only into staging, validate row/\allowbreak{}hash/\allowbreak{}provenance coverage, swap current read-only mirror atomically, and archive backup/\allowbreak{}validation reports. On red CI, stop before production/\allowbreak{}current upload/\allowbreak{}import, emit failure evidence to the phase-state controller, publish visible read-only failure history through CEX16 and link the generated repair predecessor. & Green deployment exits only if previous current snapshot remains rollbackable, public API is read-only, every row carries commit/\allowbreak{}toolchain/\allowbreak{}hash provenance and failed imports leave the old current snapshot active. Red CI exits through CEX15 repair-phase generation plus CEX16 failure-history publication, not production/\allowbreak{}current deployment.\\
\hline
\addlinespace[0.24em]
CI build-result phase-transition workflow & A GitHub Actions Lean/\allowbreak{}Lake build and declared scan pipeline completes for the active phase branch or PR. & active phase\_id;\newline current work branch/\allowbreak{}PR;\newline backup ref;\newline commit SHA;\newline GitHub Actions run/\allowbreak{}job conclusion;\newline Lean build logs;\newline scanner/\allowbreak{}artifact manifest;\newline current phase/\allowbreak{}object/\allowbreak{}proof/\allowbreak{}scratchpad registers;\newline dependency graph;\newline last green snapshot. & If CI is red, block merge and production/\allowbreak{}current deploy/\allowbreak{}import, set the attempted active phase to R, record failure evidence, publish a visible read-only failure-history snapshot, and create or activate a predecessor repair phase at the natural source layer with object-link policy, dependencies, scratchpad and proof-dossier stubs. If CI is green, record build/\allowbreak{}scanner/\allowbreak{}artifact evidence, set the active phase to G only when closeout fields are complete, update linked object/\allowbreak{}dossier/\allowbreak{}scratchpad rows, and advance exactly one active cursor to the next dependency-complete phase. & Red exit leaves no production SQL/\allowbreak{}Explorer current update and no merge, but does leave a visible failure-history snapshot and repair predecessor as the active follow-up if dependencies are known. Green exit leaves the completed phase G, evidence complete, public mirrors updated only through validated snapshot import, and exactly one next active phase selected.\\
\hline
\addlinespace[0.24em]
Local archival and rollback workflow for automated snapshots & A repository or public database snapshot becomes a release candidate or public Explorer state. & Git backup tag/\allowbreak{}branch or commit SHA, workflow artifact, SQL dump/\allowbreak{}snapshot, snapshot\_manifest.json, validation report and release timestamp;\newline for red CI, red\_failure\_snapshot.json and repair-phase link as history artifacts rather than production/\allowbreak{}current release artifacts. & Create downloadable/\allowbreak{}locally archivable git bundle or mirror-clone instruction;\newline store SQL dump or generated SQL artifact for green release candidates;\newline record manifest hashes;\newline document restore path for Git and MySQL separately. Preserve red failure snapshots in history/\allowbreak{}archive without promoting them to current production state. & A release snapshot is not considered deployment-hardened until Git state, generated artifacts and database mirror state can each be restored independently. A red failure snapshot is considered review-visible only if archived or queryable as history while last green current remains rollbackable.\\
\hline
\addlinespace[0.24em]
Red-build failure visibility workflow & A Lean/\allowbreak{}scan CI run fails for the active phase branch/\allowbreak{}PR and the failure affects phase-state transition. & commit SHA;\newline GitHub Actions run/\allowbreak{}job IDs;\newline active phase\_id;\newline failed job names;\newline sanitized log/\allowbreak{}artifact pointers;\newline phase\_state\_transition\_manifest.json;\newline generated repair-phase ID/\allowbreak{}link;\newline last green snapshot ID/\allowbreak{}hash;\newline no-current-swap evidence. & Generate red\_failure\_snapshot.json or equivalent;\newline insert/\allowbreak{}read-only expose a failure-history row/\allowbreak{}API view;\newline link the failure to the generated predecessor repair phase or localization gate;\newline preserve last green production/\allowbreak{}current snapshot;\newline scrub secrets and avoid publishing excessive logs. & The red failure is visible to reviewers and Explorer/\allowbreak{}history consumers, but no merge, production/\allowbreak{}current MySQL swap or green-snapshot publication occurs. The failure artifact is generated-secondary evidence, not a truth source.\\
\hline
\addlinespace[0.24em]
Timestamp-bounded delta-PDF generation workflow & A planning update changes a bounded set of YAML rows under one planning\_data\_timestamp/\allowbreak{}change timestamp, but no full-release or full-audit checkpoint is requested or explicitly triggered by the user. & planning\_data\_version;\newline planning\_data\_timestamp;\newline per-row Change version;\newline per-row Change timestamp;\newline changed-row section/\allowbreak{}key list;\newline generated TSV exports;\newline optional context-slice query for the same change version. & Mark every changed row with the same Change version and Change timestamp;\newline regenerate TSV/\allowbreak{}table exports;\newline emit a delta PDF containing only matching changed rows/\allowbreak{}sections and a matched-row export bundle;\newline do not compile or present the full documentation PDF as the default artifact. & The user receives the updated full ZIP plus a delta PDF. The delta PDF contains only records whose Change version/\allowbreak{}Change timestamp match the update set, with an explicit note that Lean remains primary and the delta is generated-secondary. Full PDF generation is optional, explicit or release/\allowbreak{}full-audit-gated and must not be silently substituted as the default bounded-change artifact.\\
\hline
\addlinespace[0.24em]
Workflow coherence, precedence and anti-self-sabotage audit workflow & Before adding, deleting or materially changing workflow\_policy rows;\newline before activating CI/\allowbreak{}DB/\allowbreak{}Explorer automation;\newline after any user identifies a possible workflow conflict or self-blocking rule. & Complete workflow\_policy section;\newline global\_no\_go\_guardrails;\newline active phase table;\newline object/\allowbreak{}proof/\allowbreak{}scratchpad registers;\newline CI red/\allowbreak{}green state-machine objects CEX15-CEX16;\newline database/\allowbreak{}read-only mirror policy;\newline documentation-output policy;\newline current latest\_change\_overview. & Classify every affected workflow by state transition domain (preparation, implementation, build evidence, red failure, green closeout, DB mirror, backup/\allowbreak{}rollback, documentation). Check for contradictory exits, missing predecessor dependencies, hidden red paths, false-green paths, direct-main writes, production/\allowbreak{}current swaps on red, public write leaks and full-PDF churn. Record every contradiction as a static finding and patch the earliest workflow/\allowbreak{}source layer, not downstream consumers. & The workflow set passes only if missing-evidence, red-failure, green-success and documentation-delta states remain distinct;\newline red failures are visible but never current production;\newline green closeout requires evidence and exactly one active cursor;\newline public DB/\allowbreak{}API remains read-only;\newline backup/\allowbreak{}rollback precedes mutable repository or current-mirror changes;\newline bounded planning changes emit delta PDFs by default.\\
\hline
\addlinespace[0.24em]
\end{longtable}
\arrayrulecolor{black}
\normalsize
